\documentclass[11pt]{article}
\usepackage[a4paper,margin=28mm]{geometry}
\usepackage{amsmath,amssymb,amsthm,mathtools,hyperref}
\newcommand{\ddc}{dd^c}
\newcommand{\MA}{\operatorname{MA}}
\newcommand{\capH}{\operatorname{cap}}
\newcommand{\PSH}{\operatorname{PSH}}
\newcommand{\Dzero}{\mathcal D_h^0}
\newtheorem{theorem}{Theorem}
\newtheorem{lemma}[theorem]{Lemma}
\theoremstyle{remark}\newtheorem{remark}[theorem]{Remark}
\title{Capacity domination, critical normalizers\newline and density-root tails}
\author{}\date{9 October 2026}
\begin{document}\maketitle
\noindent\textbf{Working draft.} Independent reconstruction and mathematical
review are pending. The capacity assumption below is additional to the
unrestricted root-space hypotheses in original S04.
\begin{abstract}
Under a power capacity bound on an absolutely continuous datum, we prove
a quantitative lower estimate for every dominated density's critical
normalizer. It controls the original supremum over all smooth positive
Gauduchon tests and gives order-continuous density roots. The clipped
equations have a uniformly convergent subsequence and solve the exact
equation at the critical coefficient in the bounded continuous class.
The proof separates scalar convergence from convergence of the actual
Hermitian products. On every positive Hermitian background an explicit
logarithmic reference gives a density outside every higher $L^p$ class,
with critical coefficient one and a quantitative root-tail estimate.
No positive lower bounded-potential mass or closedness is assumed.
The capacity and comparison estimates are imported with their precise
scope; no literature priority or unrestricted S04 completion is claimed.
\end{abstract}

\section{Conventions and the additional capacity condition}
Let $X$ be connected and compact of complex dimension $n\geq3$, $h$ a
smooth positive Hermitian form, and $dV$ a smooth strictly positive
probability volume. Use $\ddc=i\partial\bar\partial/(2\pi)$ and genuine
canonical psh-plus-smooth representatives. For bounded $h$-psh functions
the actual bounded product is $\MA_h(u)=(h+\ddc u)^n$. Define
\[
 \capH_h(E)=\sup_{w\in\PSH(X,h),\ 0\leq w\leq1}\MA_h(w)(E).
\]
Its total capacity is finite by local bounded-product estimates.

We use the original S04 normalizer conventions as background. For
nonzero $g\geq0$ in $L^1(dV)$, the continuous equations for
$g_j=\min(g,j)$ have positive normalizers $c_j$ decreasing to $b_h(g)$,
and normalized continuous genuine potentials $\phi_j$ with
\[
 \MA_h(\phi_j)=c_jg_jdV,\qquad \sup_X\phi_j=0.
\]
The clipped existence statement is \cite[Theorem 0.1]{KN}.
The identification of this scalar limit with the maximal genuine
subsolution coefficient is part of the original threshold background.
For real or complex $p\in L^n(dV)$ the original all-test norm is
\[
 \rho_h(p)=b_h(|p|^n)^{-1/n}
 =\sup_T\int_X|p|\left(\frac{\gamma^n}{dV}\right)^{(n-1)/n}dV,
 \qquad T=\gamma^{n-1},
\]
where the supremum retains \emph{all} smooth positive $(n-1,n-1)$ forms
$T$ with $\ddc T=0$ and $\int h\wedge T=1$. The zero-root norm is zero;
a zero threshold gives infinite norm. The space $\Dzero$ is the closure
of bounded measurable roots for $\|p\|_n+\rho_h(p)$. We use the original
equivalence between membership and vanishing root-tail norm.
These scalar/test identities are background, not conclusions of the
capacity argument.

\begin{theorem}\label{th:main}
Suppose $0\ne f\geq0$ belongs to $L^1(dV)$ and, for constants
$A,\alpha>0$, satisfies
\begin{equation}\label{eq:power}
 \int_E f\,dV\leq A\capH_h(E)^{1+\alpha}
 \quad\text{for every Borel }E\subset X.
\end{equation}
There is $C>0$, depending only on $X,h,A,\alpha$, such that every
nonzero measurable $0\leq g\leq f$ satisfies
\begin{equation}\label{eq:critical}
 b_h(g)\geq C\left(\int_Xg\,dV\right)^{-\alpha/(1+\alpha)},
 \qquad
 \rho_h(g^{1/n})\leq C^{-1/n}
       \left(\int_Xg\,dV\right)^{\alpha/(n(1+\alpha))}.
\end{equation}
In particular $f^{1/n}\in\Dzero$. There is a bounded continuous
genuine $h$-psh $\phi$, normalized by $\sup_X\phi=0$, satisfying
\begin{equation}\label{eq:exact}
 \MA_h(\phi)=b_h(f)f\,dV.
\end{equation}
The coefficient is finite and strictly positive, and the actual measure
has finite positive mass. No uniqueness assertion for homogeneous
potentials is made.
\end{theorem}

\section{Capacity estimates and critical coefficients}
We specify the two imported inequalities that will also give uniform
stability. Let $\theta>0$ be any fixed smooth Hermitian background.
For bounded genuine $a\leq0$, $-1\leq v\leq0$, fix $0<\varepsilon<1$,
put
\[
 m=\inf_X(a-(1-\varepsilon)v),\qquad
 U_\varepsilon(s)=\{a<(1-\varepsilon)v+m+s\}.
\]
The permitted increments below lie in fixed positive intervals depending
only on $\varepsilon$, dimension and the smooth torsion bounds of
$\theta$, as stated in \cite[Theorem 5.3 and Lemma 5.4]{KN}.
If $\MA_\theta(a)\leq A_0\capH_\theta^{1+\alpha}$, those statements give
\begin{align}
 s&\leq C_\alpha A_0^{1/n}
                 \capH_\theta(U_\varepsilon(s))^{\alpha/n},
                 \label{eq:lowercap}\\
 [(1-\varepsilon)t]^n\capH_\theta(U_\varepsilon(s))
  &\leq C\MA_\theta(a)
          (U_\varepsilon(s+4(1-\varepsilon)t)).\label{eq:uppercap}
\end{align}
For \eqref{eq:lowercap}, use the source admissible function
$H(x)=A_0^{-1}x^{n\alpha}$ in its displayed formula for $\kappa$.
The integral $\int_{s^{-1/n}}^\infty dx/(xH(x)^{1/n})$
equals $A_0^{1/n}s^{\alpha/n}/\alpha$, which proves the displayed
dependence, including small $A_0$. Constants in
\eqref{eq:uppercap} do not depend on the potential's oscillation.
Both source statements require a bounded competitor, and that bound
will always be checked.

The source convention is $\ddc_{\rm source}=2\ddc$.
Apply its statements to the background $2\theta$ and unchanged
potentials: the actual products and capacities are multiplied by $2^n$.
Thus the permitted intervals and constants change by fixed geometric
factors and the two displayed structural inequalities follow exactly.
This conversion also applies to the clipped existence and weighted
CLN estimates. No closedness or positive lower mass is introduced.

\begin{proof}[Proof of the coefficient estimate]
Let $g_j=\min(g,j)$ and use the clipped potentials above. Their measures
are bounded by $c_jA\capH_h^{1+\alpha}$. Set $\varepsilon=1/2$,
use the competitor $v=0$, and choose fixed positive $s,t$ in \emph{both}
permitted intervals. Then \eqref{eq:lowercap} and
\eqref{eq:uppercap}, the latter bounded by total mass, give
\[
 s\leq C(c_jA)^{1/n}\capH_h\{\phi_j<\inf\phi_j+s\}^{\alpha/n},
 \qquad
 \capH_h\{\phi_j<\inf\phi_j+s\}
    \leq Ct^{-n}c_j\int_Xg_jdV.
\]
Combining yields
$s^n\leq C A c_j^{1+\alpha}(\int g_jdV)^\alpha$, and hence
the first inequality in \eqref{eq:critical} upon taking the original
scalar limit $c_j\downarrow b_h(g)$. The original threshold/test
identity gives the second inequality over the full test family.

Take $g=f\mathbf1_{\{f>K\}}$. Its mass tends to zero. The actual root
tail is
\[
 g^{1/n}=f^{1/n}\mathbf1_{\{f^{1/n}>K^{1/n}\}};
\]
both its $L^n$ norm and its $\rho_h$ norm tend to zero.
This proves $f^{1/n}\in\Dzero$ under the original norm.
\end{proof}

\section{Uniform compactness and the actual exact equation}
For $f_j=\min(f,j)$, let $c_j$ and $\phi_j$ be the clipped data.
The coefficient estimate gives $b_h(f)>0$, while $c_j\leq c_1<\infty$.
All measures therefore satisfy one common power bound
$Ac_1\capH_h^{1+\alpha}$. Applying \eqref{eq:lowercap} with $v=0$
at one fixed $s>0$ gives
\[
 \capH_h\{\phi_j<\inf\phi_j+s\}\geq c_0>0.
\]
Weighted CLN \cite[Proposition 1.1]{Ngu}, for normalized genuine
$h$-psh $\phi_j$ and every $0\leq w\leq1$, gives
$\int|\phi_j|\MA_h(w)\leq C$. Consequently
\[
 \capH_h\{\phi_j<-L\}\leq C/L.
\]
If $-\inf\phi_j>s$, the former sublevel lies in the latter with
$L=-\inf\phi_j-s$; thus these two bounds prove
\begin{equation}\label{eq:bound}
 -M\leq\phi_j\leq0
\end{equation}
for one $M\geq1$. This does not use a uniform $L^p$ norm of $f_j$.

\begin{lemma}\label{lem:stability}
Let $u_j\in\PSH(X,\theta)\cap C(X)$ lie in $[-1,0]$, with
$\MA_\theta(u_j)\leq A_0\capH_\theta^{1+\alpha}$ uniformly.
If
\[
 \int_X|u_j-u_k|\MA_\theta(u_j)\longrightarrow0
 \quad\text{as }j,k\longrightarrow\infty,
\]
then the sequence is uniformly Cauchy.
\end{lemma}
\begin{proof}
Fix $d>0$ and suppose $\inf(u_j-u_k)\leq-d$.
Choose $0<\varepsilon<\min(d/8,1/4)$ and positive $s,t$ in both
permitted intervals such that $s+4t+\varepsilon<d/2$.
Use $a=u_j$, $v=u_k$ above. Since $v\leq0$,
$m=\inf(a-(1-\varepsilon)v)\leq-d$. As $v\geq-1$,
\[
 U_\varepsilon(s+4(1-\varepsilon)t)
        \subset\{a-v<-d/2\}.
\]
Equation \eqref{eq:uppercap} therefore bounds
$\capH_\theta(U_\varepsilon(s))$ by
$C_{\varepsilon,t}d^{-1}\int|a-v|\MA_\theta(a)$,
which tends to zero. Equation \eqref{eq:lowercap} contradicts this
for the fixed $s>0$. Exchanging $j$ and $k$ proves both directions.
All parameters are fixed for each $d$, before sending the indices
to infinity.
\end{proof}

\begin{proof}[Proof of the exact equation]
Normalized genuine $h$-psh functions have a uniform $L^1(dV)$ bound
\cite[Proposition 2.1]{DK}. Add fixed smooth psh corrections on a finite
chart covering and apply local $L^1$ compactness
\cite[Chapter I, Propositions 4.21 and 5.9]{Dem}. This gives global
$L^1(dV)$ compactness for the sequence.
Extract a subsequence $\phi_j$ converging in $L^1(dV)$ and $dV$-almost
everywhere. In view of \eqref{eq:bound} and the finite measure $fdV$,
dominated convergence gives the two-index limit
\[
 \int|\phi_j-\phi_k|\MA_h(\phi_j)
   \leq c_1\int|\phi_j-\phi_k|f\,dV\longrightarrow0.
\]
This uses absolute continuity of the fixed datum; it does not multiply
an arbitrary fixed Borel function by a weakly convergent measure.

Put $\theta=h/M$ and $u_j=\phi_j/M$. They lie in $[-1,0]$ and
$\MA_\theta(u_j)=M^{-n}\MA_h(\phi_j)$.
For each $h$-capacity competitor $w$, the function $w/M$ is a
$\theta$-competitor, so
$\capH_h(E)\leq M^n\capH_\theta(E)$.
Thus the scaled products obey one common power capacity bound,
and their two-index errors tend to zero. Lemma~\ref{lem:stability}
makes the subsequence uniformly Cauchy. Its limit $\phi$ is continuous,
genuinely $h$-psh by distributional positivity, and has supremum zero.

On every chart choose a fixed smooth strictly psh $q$ with
$\ddc q\geq h$; $\phi_j+q$ are uniformly bounded psh. Expand
\[
 \MA_h(\phi_j)
 =\sum_{\ell=0}^n\binom n\ell
      [\ddc(\phi_j+q)]^\ell\wedge(h-\ddc q)^{n-\ell}.
\]
For completeness put $S_j=\ddc(\phi_j+q)$ and $S=\ddc(\phi+q)$.
For each $\ell\geq1$, bounded multilinearity and the closed-product
recursion \cite[Section 2]{BT} give
\[
 S_j^\ell-S^\ell=\ddc\left[(\phi_j-\phi)
          \sum_{a=0}^{\ell-1}S_j^a\wedge S^{\ell-1-a}\right].
\]
Pairing with a compactly supported smooth test form bounds this by
$C\|\phi_j-\phi\|_\infty$: local CLN bounds all the mixed masses
uniformly \cite[Theorem 2.10]{BT}. Thus every closed local product
converges weakly, with all smooth correction factors retained.
This is the actual bounded Hermitian product of $\phi$.
Meanwhile $c_jf_j\to b_h(f)f$ in $L^1(dV)$, since $c_j\leq c_1$,
$c_j\to b_h(f)$ and $f_j\to f$ in $L^1$.
The two limits prove \eqref{eq:exact}, including its exact coefficient
and finite prescribed measure in the bounded continuous class.
\end{proof}

\section{An explicit density beyond every higher integrability class}
\begin{theorem}\label{th:example}
On every $X,h,dV$ as above and for each $\alpha>0$, there is a bounded
continuous genuine $h$-psh $u$, with $h+\ddc u\geq h/2$, whose actual
product is $fdV$ with $f\in L^1(dV)$ and $f\notin L^p(dV)$ for
every $p>1$. This density satisfies \eqref{eq:power}, has
$b_h(f)=1$, and obeys, for sufficiently large $K$,
\[
 \rho_h\bigl(f^{1/n}\mathbf1_{\{f>K\}}\bigr)
          \leq C(\log K)^{-\alpha}.
\]
This is also an example on any given background with $v_-(h)=0$;
no assertion that all backgrounds have zero lower mass is made.
\end{theorem}
\begin{proof}
Fix a small holomorphic ball about $p$, with $R=|z|<1$, and a smooth
cutoff $0\leq\chi\leq1$ equal to one on a smaller ball. Put
\[
 q(s)=(1-s)^{-\alpha},\quad s=\log R^2,
 \qquad \psi=\chi q(\log R^2)+1-\chi.
\]
Set $q(-\infty)=0$, and extend $\psi=1$ outside the cutoff ball.
The function is bounded continuous; in the inner ball,
$q'=\alpha(1-s)^{-\alpha-1}>0$ and
$q''=\alpha(\alpha+1)(1-s)^{-\alpha-2}>0$, so it is psh.
The transition is smooth with $\ddc\psi\geq-K_0h$ for one fixed
$K_0$. Choose $\delta>0$ small enough and set
$u=\delta\psi-\sup_X\delta\psi$; then
$h+\ddc u\geq h/2$ everywhere as actual currents.

The bounded product has no mass at $p$. Indeed, after a fixed smooth
local correction the positive mixed measures are bounded by a constant
times the top product of a bounded psh function plus a smooth strictly
psh function. Normalize that sum to lie in $[0,1]$ on the chart. Its
top measure is bounded by its local capacity, and the pluripolar point
has capacity zero \cite[Theorem 6.9 and the capacity definition]{BT}.
The finite smooth-correction expansion therefore has no atom there.
Off $p$ the reference is smooth, so the
actual product is absolutely continuous globally. Set $fdV=\MA_h(u)$.
Positivity and the strict lower bound show $f\geq c>0$ almost everywhere
for a fixed smooth-volume constant; in particular its measure has full
support. In the inner ball the pure top term has Euclidean density
\begin{equation}\label{eq:radial}
 (\ddc u)^n=
 \frac{\delta^n n!\alpha^n(\alpha+1)}{\pi^n}
 R^{-2n}(1-\log R^2)^{-n(\alpha+1)-1}\,d\lambda.
\end{equation}
This follows from the tangential eigenvalues $q'/R^2$ and the radial
eigenvalue $q''/R^2$. Integrating the sphere gives exactly
\[
 \int_{B_r}(\ddc u)^n
     =\delta^n\alpha^n(1-\log r^2)^{-n(\alpha+1)}.
\]
All lower degree $k<n$ mixed densities are bounded by
$C R^{-2k}(1-\log R^2)^{-k(\alpha+1)}$ and have ball mass $O(r^2)$.
The inner expansion with the smooth positive $h$ has positive terms.
It follows that $f\in L^1$ and, from the lower bound
\eqref{eq:radial}, that $f\notin L^p$ for any $p>1$.
For sufficiently small $r$ we also obtain
\begin{equation}\label{eq:ball}
 \int_{B_r}f\,dV\leq C[(1-\log r^2)^{-n(\alpha+1)}+r^2],
 \qquad \mathop{\rm ess\,sup}_{X\setminus B_r}f\leq C r^{-2n}.
\end{equation}

The volume-capacity bound \cite[Proposition 5.1]{KN} gives
\[
 dV(E)\leq C\exp[-a\capH_h(E)^{-1/n}].
\]
Changing $h^n$ to the fixed comparable smooth $dV$ changes constants
only. The primary antecedent \cite[Corollary 2.4]{DK}, applied to the
constant density, gives the same exponential estimate on compacts;
inner regularity of smooth volume and capacity monotonicity give it
for all Borel sets. For $s=\capH_h(E)>0$ small, choose
$r=\exp[-a s^{-1/n}/(4n)]$, inside the permitted ball.
Split $E$ inside and outside $B_r$ and use \eqref{eq:ball}:
\[
 \int_Ef\,dV\leq C s^{1+\alpha}
                +C\exp[-a s^{-1/n}/2]
          \leq C' s^{1+\alpha}.
\]
The exponentially small $r^2$ term is absorbed by the power term.
For larger $s$ enlarge the constant using the finite total mass.
For $s=0$ the volume bound implies $dV(E)=0$, so the datum vanishes
there. This verifies \eqref{eq:power} for every Borel set.

Theorem~\ref{th:main} gives a continuous normalized $\phi$ satisfying
$\MA_h(\phi)=b_h(f)fdV$. The explicit $u$ is already a genuine
subsolution with coefficient one, so the original maximal-coefficient
background gives $b_h(f)\geq1$. If $b_h(f)>1$, choose small
$\varepsilon>0$ so that $(1-\varepsilon)^n b_h(f)>1$.
The bounded positive expansion gives
$\MA_h((1-\varepsilon)\phi)\geq(1-\varepsilon)^n b_h(f)fdV$.
Modified comparison \cite[Theorem 2.3]{KN}, for $u$ and
$(1-\varepsilon)\phi$ on
\[
 U=\{u<(1-\varepsilon)\phi+
                 \inf_X(u-(1-\varepsilon)\phi)+s\},
\]
with $s>0$ sufficiently small, gives
\[
 (1-\varepsilon)^n b_h(f)\int_U f\,dV
       \leq(1+C_\varepsilon s)\int_U f\,dV.
\]
The permitted increment interval is positive for fixed $\varepsilon$;
choose $s$ in it with $1+C_\varepsilon s<(1-\varepsilon)^n b_h(f)$.
The set $U$ is open and nonempty by continuity and attainment of the
minimum. Its measure is positive because $f\geq c>0$.
This contradiction proves $b_h(f)=1$.

Finally take $r_K=K^{-1/(4n)}$. For large $K$, \eqref{eq:ball}
gives $\mathop{\rm ess\,sup}_{X\setminus B_{r_K}}f\leq C K^{1/2}<K$,
so $\{f>K\}\subset B_{r_K}$ up to a null set. Hence
\[
 \int_{\{f>K\}}f\,dV
    \leq C(\log K)^{-n(\alpha+1)}+CK^{-1/(2n)}
    \leq C'(\log K)^{-n(\alpha+1)}.
\]
Applying \eqref{eq:critical} to $f\mathbf1_{\{f>K\}}$ gives the
asserted root-tail rate without an implicit threshold inversion.
\end{proof}

\begin{remark}
The logarithmic reference is continuous and bounded but is not H\"older
at $p$. The example therefore goes beyond higher $L^p$ integrability
without invoking a H\"older-reference range theorem. The analytic
capacity estimates and their historical range consequences are prior
work. The explicit estimate on the original critical normalizer and the
actual product argument are stated with their additional capacity
hypothesis. Original S04 asks whether \emph{every} positive-threshold
root has uniform tails, or whether uniform tails alone give the exact
critical equation. This paper proves neither universal assertion and
does not supply a counterexample to either one.
\end{remark}

\begin{thebibliography}{9}
\bibitem{KN} S. Ko\l odziej and N. C. Nguyen,
\emph{Weak solutions to the complex Monge--Amp\`ere equation on Hermitian
manifolds}, \href{https://arxiv.org/abs/1312.5491v2}{arXiv:1312.5491v2},
Theorem 0.1, Theorem 2.3, Proposition 5.1, Theorem 5.3 and Lemma 5.4.
\bibitem{Ngu} N. C. Nguyen,
\emph{The complex Monge--Amp\`ere type equation on compact Hermitian
manifolds and applications},
\href{https://arxiv.org/abs/1501.00891v1}{arXiv:1501.00891v1}, Proposition 1.1.
\bibitem{DK} S. Dinew and S. Ko\l odziej,
\emph{Pluripotential estimates on compact Hermitian manifolds},
\href{https://arxiv.org/abs/0910.3937v1}{arXiv:0910.3937v1},
Proposition 2.1, Corollary 2.4 and Proposition 2.5.
\bibitem{BT} E. Bedford and B. A. Taylor,
\emph{A new capacity for plurisubharmonic functions}, Acta Math.
149 (1982), 1--40, Sections 2, 6 and 7.
\bibitem{Dem} J.-P. Demailly,
\emph{Complex Analytic and Differential Geometry}, author-hosted
manuscript, version of 21 June 2012, Chapter I, Propositions 4.21 and 5.9.
\end{thebibliography}
\end{document}
